\section{From row comparisons to streaming lower bounds}
\label{sec:r-streaming}

To apply the row comparisons to the learner in
Definition~\ref{def:r-streaming-model}, we first replace
completed-measurable rules by Borel versions in the selected
uniform-prior experiment.  We then account for the stopping index
and quantify the information that an accurate output must retain
after independent auxiliary rows are revealed.  These common
reductions leave each comparison its own block estimate, row
dimensions, and terminal residual-radius bound.

\subsection{Borel rules after the shared seed is fixed}
\label{subsec:r-borel-reduction}

Lemma~\ref{lem:r-seed-reduction} is applied to the original jointly
measurable experiment.  It supplies a fixed seed for which
\eqref{eq:r-seed-factorization} holds and the required uniform-prior
success is retained.  The following reduction is performed only
after that choice.  It does not assert preservation of behavior at
every individual signal.

\begin{lemma}[Borel versions of the finite-state rules]
\label{lem:r-borel-kernels}
Fix a seed satisfying \eqref{eq:r-seed-factorization}, a dimension
\(d\ge2\), and a deterministic finite horizon \(T\).  Retain the
conditional law of \(U_0\) at that seed.  Suppose each transition
and stopping probability at a positive sample index is measurable
in the completion of the Borel sigma field on
\(\R^d\times\R\) under
\[
 \lambda(dx,dy)=\gamma_1(dx)\,dy .
\]
There are everywhere defined Borel transition and stopping kernels
with the same state sets and terminal output kernels whose
fixed-seed uniform-prior joint law of the signal, state path,
stopping index, and output is the same as before.  Their block
routing probabilities are Borel functions of the entering state
and the entire block data, including at states and data that have
zero probability in the original run.
\end{lemma}

\begin{proof}
For each state and positive index, combine the next-state and
stopping choices into one finite destination set, tagging a
destination as continuing or terminal.  The kernel is a finite
vector of completed-measurable nonnegative functions of \((x,y)\)
whose sum is one.  Every completed-measurable real function has a
Borel version: approximate it by simple functions and replace their
level sets by Borel sets modulo null sets.  Choose such a version
for every coordinate.  The union of the exceptional sets is
contained in a Borel \(\lambda\)-null set.  On the Borel set where
the chosen vector is not a probability vector, replace it by a
fixed point mass.  That set is also null.  Since the horizon and all
state sets are finite, a single Borel \(\lambda\)-null set \(Z\)
contains every sample-pair modification.  A stopping choice at
index zero depends only on the initial state and the fixed seed,
with fresh randomness, and has no sample argument to replace.

In the fixed-seed law, \eqref{eq:r-seed-factorization} leaves
\(S\sim\sigma\) independent of the pre-generated Gaussian rows.
For a nonzero row \(x\), the image of \(\sigma\) under
\(s\mapsto\langle x,s\rangle\) has a Lebesgue density when \(d\ge2\).
This follows from the one-coordinate sphere formula after rotation
and multiplication by \(\|x\|\).  The zero row has Gaussian
probability zero.  Hence every marginal sample pair
\((X_t,\langle X_t,S\rangle)\) is absolutely continuous with
respect to \(\lambda\).  In particular,
\[
 \Prob\{(X_t,\langle X_t,S\rangle)\in Z\}=0.
\]
For any state \(u\), even one selected using previous observations,
the probability that the state before time \(t\) is \(u\) and the
current pair lies in \(Z\) is bounded by this same zero probability.
A finite union over the pre-generated indices therefore shows that
no modified input is visited almost surely.

Couple the original and Borel finite kernels with the same fresh
uniform random variable at each step.  Induction makes their state
paths and stopping choices identical almost surely.  At identical
terminal state-index pairs their output kernels agree, so the
outputs may also be coupled identically.  This proves the claimed
joint-law equality.  For any specified entering state, finite
composition of the everywhere Borel kernels gives Borel block
routing probabilities on every block-data vector.  Such a routing
function is a continuation from the specified state; it need not
be a conditional law given that the state was reached.
\end{proof}

The order of these reductions matters.  Original joint measurability
makes the conditional success function of the shared seed
measurable, so it can be averaged in Lemma~\ref{lem:r-seed-reduction}.
Lemma~\ref{lem:r-borel-kernels} then chooses versions separately for
the selected seed.  No jointly measurable selection of replacement
versions across seed values is needed.  In particular the global
Borel extensions used in the finite fiber comparison are fixed
before its singular comparison measure is constructed; no
real-law null set is transferred to a perpendicular Gaussian law.

In each fixed-seed experiment below, unused future rows are still
independent of the joint vector consisting of the signal and the
current state.  This follows from the pre-generated product row law,
\eqref{eq:r-seed-factorization}, and the fresh transition
randomness in the model.  If stopping is allowed at index zero,
let \(\overline U_0\) denote the initial padded state, including
that status.  It is generated from \(U_0\) and fresh time-zero
randomness, so
\[
 \overline U_0\ \perp\!\!\!\perp\
       (S,(X_t)_{t\ge1})
\]
in the fixed-seed law.  Whenever an auxiliary matrix is introduced,
it is sampled independently of the pair consisting of \(S\) and
the state in question.  These joint independence statements justify
the zero initial potentials and the fresh-block comparisons.

\subsection{Common stopping and terminal estimates}
\label{subsec:r-stopping-reduction}

The five comparisons use different information potentials.  Their
application to a learner shares two estimates.  First we account for
the stopping index before any analytic comparison is invoked.  We then
bound the information that success requires after auxiliary rows have
been revealed.

\begin{lemma}[Padding, output capacity, and unrestricted data]
\label{lem:r-stopping-reduction}
Consider a fixed-seed Borel learner from
Lemmas~\ref{lem:r-seed-reduction} and \ref{lem:r-borel-kernels}.
Suppose its accuracy satisfies \(0<\epsilon\le1/10\), its
persistent states have at most \(2^M\) values, and its deterministic
finite horizon is \(T\).  Write
\(L=\log(1/\epsilon)\).  It admits an absorbing extension,
including to the end of a longer final block, with at most
\[
 N=(T+2)2^M
\]
values at each layer.  The terminal output is a kernel of the extended
state alone.  The initial extended state is independent of the joint
signal and row sequence.  At every deterministic block boundary,
unused future rows are independent of the pair consisting of the
signal and the entering extended state.  The block transition is a
Borel kernel of that state and the exact block data.

If its uniform-prior angular success probability is \(p\), then
\begin{equation}\label{eq:r-common-output-capacity}
 p\le (T+1)2^M\epsilon^{d-1}.
\end{equation}
For a sequence of such learners with \(p\ge3/5\) and
\(M(d)=o(d^2)\), on the branch \(T\le dL\),
\begin{equation}\label{eq:r-common-clock}
 dL=O(M+1),\qquad L=o(d),\qquad
 \log N=M\log2+\log(T+2)=o(d^2).
\end{equation}
All these estimates are uniform over the allowed accuracies.
For all sufficiently large \(d\), even a learner given all its exact
observations must have \(T>d/4\) to achieve success at least \(3/5\).
\end{lemma}

\begin{proof}
After stopping, retain the terminal state together with its stopping
index in an absorbing state.  Defer the output until the end of the
extension and apply its original state-index kernel.  Active states
and all stopped pairs number at most \((T+2)2^M\).  At index zero the
stopping status is generated from the data-independent initialization
and fresh randomness.  The independence statements then follow from
the pre-generated product row law.  Finite composition of the Borel
rules gives the asserted block kernel, including on inputs never
visited in the original run.

There are at most \((T+1)2^M\) terminal state-index labels.
For each label, its output kernel has no further dependence on the
signal.  In the success integral for that label, bound the conditional
probability of reaching it by one.  Lemma~\ref{lem:r-caps}, averaged
over the output kernel, bounds the remaining integral by
\(\epsilon^{d-1}\).  Summing proves
\eqref{eq:r-common-output-capacity}, including for randomized output.
When \(p\ge3/5\), taking logarithms gives
\[
 (d-1)L\le M\log2+\log(T+1)+\log(5/3).
\]
For sufficiently large \(d\), uniformly for \(L\ge\log10\),
\((d-1)L\ge3dL/4\) and \(\log(dL+1)\le dL/4\).
On \(T\le dL\), subtraction gives \(dL=O(M+1)\).
Consequently \(L=o(d)\), \(T=o(d^2)\), and
\(\log N=o(d^2)\).  This is a conclusion only on the stated
branch; no restriction on the accuracy sequence has been imposed.

For the unrestricted-data bound, suppose \(T\le d/4\) and grant
the output all \(T\) pre-generated rows and labels, including unused
pairs after a stop.  Given those data, the remaining signal is uniform
on a sphere of radius \(R\) in a \((d-T)\)-dimensional affine space.
Indeed the row-space component is determined, while rotational
invariance leaves a uniform direction in the kernel.  More generally,
if \(U\) is the span of \(t<d\) independent Gaussian rows, isotropy
and independence give
\begin{equation}\label{rlevels:residual-mean}
 \mathbb E\|P_US\|^2=t/d.
\end{equation}
Taking \(t=T\) gives \(\mathbb E(1-R^2)=T/d\), so
\[
 \Pr\{R<1/2\}\le\frac{4T}{3d}\le\frac13.
\]
On the complementary event, an output ball of radius \(\epsilon\)
intersecting that residual sphere is contained in a ball of radius
\(2\epsilon\) centered at a residual-sphere point.  Rescaling by
\(R\ge1/2\) and applying Lemma~\ref{lem:r-caps} bounds conditional
success by \((4\epsilon)^{d-T-1}\).  The learner's remaining
randomness has no further signal information after all data are given,
so the same estimate holds after averaging its output.  Thus success
is at most
\[
 \frac13+(2/5)^{3d/4-1}<\frac35
\]
for large \(d\), a contradiction.  This argument includes \(T=0\),
when the residual sphere is the original unit sphere.
\end{proof}

\begin{lemma}[Residual information in a successful state]
\label{lem:r-residual-terminal}
Let \(S\sim\sigma\), let \(V\) be a finite state jointly
distributed with \(S\), and let \(G\) have \(0\le k<d-1\)
independent standard Gaussian rows, independently of \((S,V)\).
Suppose a unit-vector output is drawn from a kernel of \(V\) alone
and has angular success probability at least \(p\) at accuracy
\(\epsilon>0\).  Let \(R\) be the radius of the conditional
sphere of \(S\) given \((G,GS)\).  If
\[
 0<\rho\le1,\qquad 2\epsilon<\rho,\qquad
 \Pr\{R<\rho\}\le\eta<p,
\]
then
\begin{equation}\label{eq:r-residual-terminal}
 I(S;V\mid G,GS)
 \ge(p-\eta)(d-k-1)\log\frac{\rho}{2\epsilon}-\log2.
\end{equation}
For \(\rho<1\), one available bad-radius bound is
\begin{equation}\label{eq:r-residual-markov}
 \Pr\{R<\rho\}\le\frac{k}{d(1-\rho^2)}.
\end{equation}
\end{lemma}

\begin{proof}
Let \(Q\) preserve the side-data marginal \((G,GS)\) and the
conditional marginals of \(S,V\) given those data, while making
\(S\) and \(V\) conditionally independent.  Apply the same
state-output kernel under both laws.  Their relative entropy,
including the output, is \(I(S;V\mid G,GS)\).

The event of angular success and \(R\ge\rho\) has actual
probability at least \(p-\eta\).  Under \(Q\), condition also
on the output.  The residual signal remains uniform on its sphere
in the kernel of \(G\), independently of the output.  Angular
success implies Euclidean error at most \(\epsilon\).  If the
corresponding output ball intersects the residual sphere, centering
it at one point of the intersection enlarges its radius to at most
\(2\epsilon\).  Rescaling the residual sphere to unit radius
and applying Lemma~\ref{lem:r-caps} bounds the conditional success
probability by \((2\epsilon/\rho)^{d-k-1}\).

For any event with actual probability \(a\) and reference
probability at most \(b\), where \(0<b\le1\), binary relative
entropy and data processing give
\begin{equation}\label{rlevels:binary-kl}
 D(P\Vert Q)\ge a\log(1/b)-\log2.
\end{equation}
Indeed the binary entropy of \(a\) is at most \(\log2\).
Applying this inequality to the preceding event proves
\eqref{eq:r-residual-terminal}.  Finally
\eqref{rlevels:residual-mean} gives
\(\mathbb E(1-R^2)=k/d\); Markov's inequality proves
\eqref{eq:r-residual-markov}.
\end{proof}

We now apply the comparisons to these extended states.  The
critical-radius, density-level, and hybrid comparisons use an entropy
bound for the outgoing state; the fiber comparison uses its number of values.  The
same width bound supplies both hypotheses.  Each route retains its own
auxiliary dimension and terminal residual-sphere estimate.

\subsection{The critical-radius comparison}
\label{subsec:r-critical-application}

\begin{theorem}[Streaming consequence of the critical radius]
\label{thm:r-critical-streaming}
Let \(M(d)\) be a nonnegative integer sequence with \(M(d)=o(d^2)\),
and let \(0<\epsilon(d)\le1/10\).  For each \(d\), consider a learner
in Definition~\ref{def:r-streaming-model} with deterministic finite
horizon \(T(d)\).  If its uniform-sphere angular success probability,
including all learner randomness, is at least \(3/5\), then
\[
 T(d)\ge c\,d\log\frac1{\epsilon(d)}
\]
for an absolute \(c>0\) and all sufficiently large \(d\).  The
eventual dimension threshold may depend on the memory sequence.
\end{theorem}

\begin{proof}
Use Lemma~\ref{lem:r-seed-reduction} with \(p_0=3/5\), and then
Lemma~\ref{lem:r-borel-kernels}.  The resulting fixed-seed
uniform-prior experiment retains success at least \(3/5\), the
conditional initial-state law, the same state bound, and the same
horizon.  Transition and terminal output randomness remain fresh.

Put \(L=\log(1/\epsilon)\ge\log10\).  The branch
\(T\ge dL\) already has the desired order.  On \(T<dL\),
Lemma~\ref{lem:r-stopping-reduction} gives padded states of entropy
at most \(H_*=M\log2+\log(T+2)\le d^2\) eventually.  In
particular its output-capacity estimate specializes to
\begin{equation}\label{eq:r-critical-output-capacity}
 \frac35\le(T+1)2^M\epsilon^{d-1}.
\end{equation}
Use \(B=\lceil T/m\rceil\) full padded blocks, where
\(m=\lfloor d/10\rfloor\) as in
\eqref{eq:r-critical-parameters}.

For a boundary state \(U\), define
\[
 J(U)=I(S;U\mid G,GS)
\]
with an independent \(\ell\)-row matrix \(G\),
\(\ell=\lfloor d/2\rfloor\), sampled independently of the pair
\((S,U)\).  In a block, write \(U,W\) for the entering and exiting
states and \(A\) for its \(m\) rows.  The matrix \(A\) is independent
of \((S,U)\).  Append an independent \(r=\ell-m\) row matrix \(C\)
and set \(G=(A;C)\).  In the aligned experiment \((G,GS)\) reveals
the entire exact block data \((A,AS)\).  Conditional on these data
and \(U\), the block transition has no further dependence on \(S\).
Conditional data processing therefore gives
\[
 I_1(S;W\mid G,GS)
 \le I_1(S;U\mid G,GS)=J(U).
\]
The equality uses independence of \(G\) from the pair \((S,U)\).
Theorem~\ref{thm:r-critical-comparison}, applied to the outgoing
state with \(H(W)\le H_*\le d^2\), yields
\[
 J(W)\le J(U)+Kd.
\]
The initial padded state \(\overline U_0\) is independent of \(S\),
and its auxiliary matrix is independent of
\((S,\overline U_0)\), so \(J(\overline U_0)=0\).  Iteration gives
for the final padded state \(W_f\)
\begin{equation}\label{eq:r-critical-iteration}
 J(W_f)\le Kd\left\lceil\frac{T}{m}\right\rceil .
\end{equation}

We next lower-bound that final information.  Conditional on
\((G,GS)\), the row-space projection of \(S\) is determined and
the residual direction is uniform in \(\ker G\).  If
\(a=d-\ell\), its radius \(R\) has the marginal representation
\[
 R^2\stackrel{\mathrm{law}}=\frac{X}{X+Y},
 \qquad X\sim\chi_a^2,\quad Y\sim\chi_\ell^2
\]
with independent \(X,Y\).  This follows by normalizing a standard
Gaussian vector.  Since \(a\ge d/2\), the events
\(X\ge0.75a\) and \(X+Y\le1.4d\) imply \(R\ge1/2\).
Chebyshev's inequality, using variances \(2a\) and \(2d\), gives
\begin{equation}\label{eq:r-critical-radius-tail}
 \Prob\{R<1/2\}
 \le\frac{32}{a}+\frac{25}{2d}=O(1/d).
\end{equation}

For sufficiently large \(d\), the radius tail in
\eqref{eq:r-critical-radius-tail} leaves success probability at
least \(1/2\) on \(R\ge1/2\).  Apply
Lemma~\ref{lem:r-residual-terminal} with \(\rho=1/2\) and
that tail bound.  It gives
\begin{equation}\label{eq:r-critical-required-information}
 J(W_f)\ge\frac{a-1}{2}\bigl(L-\log4\bigr)-\log2
          \ge c_0dL
\end{equation}
for an absolute \(c_0>0\) and all sufficiently large \(d\).
This is uniform over \(L\ge\log10\), because
\[
 L-\log4\ge
 \left(1-\frac{\log4}{\log10}\right)L>0 .
\]

The common stopping reduction gives \(T>d/4\ge m\), so
\(\lceil T/m\rceil\le2T/m\).  Equation
\eqref{eq:r-critical-iteration} is therefore an \(O(T)\) upper
bound.  Comparing it with
\eqref{eq:r-critical-required-information} proves \(T\ge cdL\)
uniformly over the stated accuracies.
\end{proof}

\subsection{The density-level comparisons}
\label{rlevels:endpoints}

We apply both comparisons to Definition~\ref{def:r-streaming-model}.
Assume that the original uniform-prior experiment has angular success
probability at least \(2/3\), that \(0<\epsilon\le1/10\), and that
\(M=M(d)=o(d^2)\); \(T\) is its deterministic finite horizon.  Put
\(L=\log(1/\epsilon)\ge\log10\).  A worst-signal success premise of
\(2/3\) first implies this uniform-prior premise by integration.

Apply Lemma~\ref{lem:r-seed-reduction} to the original jointly
measurable experiment with \(p_0=2/3\).  Fix a seed \(\xi\) for which
\eqref{eq:r-seed-factorization} holds and conditional success is at
least \(2/3\), retain \(\mathcal L(U_0\mid\Xi=\xi)\), and then apply
Lemma~\ref{lem:r-borel-kernels}.  These reductions preserve the fixed-seed
uniform-prior joint law, the state bound, and the horizon.  All learner
statements in this subsection concern that experiment.  We prove
\(T\ge c dL\) for an absolute \(c>0\) in all sufficiently large
dimensions; the eventual threshold may depend on the memory sequence.

For completeness, \eqref{rlevels:binary-kl} applied to success
under the product signal/output law and Lemma~\ref{lem:r-caps}
give the unconditional bound
\begin{equation}\label{rlevels:output-information}
 I(S;\widehat S)\ge\tfrac23(d-1)L-\log2.
\end{equation}

Restrict to \(T<dL\).  Lemma~\ref{lem:r-stopping-reduction}
provides absorbing padded states, a state-only terminal output,
data-independent initial state \(\overline U_0\), and fresh future
blocks.  Its clock bound implies, in particular,
\begin{equation}\label{rlevels:clock}
 L\le C\frac{M+1+\log d}{d},\qquad T=O(M+d),\qquad
 H_*:=M\log2+\log(T+2)=o(d^2).
\end{equation}
The lemma also gives
\begin{equation}\label{rlevels:linear}
 T>d/4.
\end{equation}

Both comparisons below use this padded learner.  Their auxiliary
matrices and block lengths remain distinct.

For the one-level comparison, return to
\eqref{rlevels:one-parameters} and use
\(B_{\mathrm{one}}=\lceil T/k\rceil\) padded blocks.  In a block let \(Z\)
and \(U\) be the entering and exiting states.  The next matrix
\(X\) is independent of the pair \((S,Z)\).  Draw \(B\) independently
of \((S,Z,X,U)\), and put \(G=(X;B)\) as in Proposition~\ref{rlevels:one-comparison}.
Conditional on \((G,GS)\), the entire block data \((X,XS)\)
are known.  The transition kernel from \(Z\) to \(U\) then has no
additional dependence on \(S\).  Conditional data processing gives
\[
 \widetilde{\mathcal D}(U)
 \le I_{\mathrm{act}}(S;Z\mid G,GS)=\mathcal D(Z),
\]
where the equality uses independence of \(G\) from \((S,Z)\).
By \eqref{rlevels:clock}, \(H(U)\le H_*\le d^2\) eventually,
so Proposition~\ref{rlevels:one-comparison} costs at most \(Cd\)
per block.  The initial boundary \(\overline U_0\) is independent of
\(S\), and its auxiliary \(G\) is independent of the pair
\((S,\overline U_0)\); hence \(\mathcal D(\overline U_0)=0\), and
\begin{equation}\label{rlevels:one-recurrence}
 \mathcal D(U_f)\le CdB_{\mathrm{one}}
\end{equation}
for the final padded state.

Take the independent \(m\)-row matrix in the definition of
\(\mathcal D\).  Its residual radius \(R\) satisfies
\[
 \Pr\{R<1/2\}\le \frac{m/d}{3/4}\le\frac13.
\]
Lemma~\ref{lem:r-residual-terminal}, with
\(p=2/3\), \(\rho=1/2\), and \(\eta=1/3\), now gives
\begin{equation}\label{rlevels:one-endpoint}
 \mathcal D(U_f)
 \ge\frac{d-m-1}{3}(L-\log4)-\log2
 \ge c_0dL
\end{equation}
for an absolute \(c_0>0\) in sufficiently large dimension.
The final inequality is uniform for \(L\ge\log10>\log4\).
Since \(T>d/4\ge k\), the common raw-data bound makes
\(\lceil T/k\rceil\le2T/k\).  Hence
\eqref{rlevels:one-recurrence} is \(O(T)\), and comparison with
\eqref{rlevels:one-endpoint} proves \(T\ge cdL\).

For the two-label comparison, use
\eqref{rlevels:two-parameters} and
\(B_{\mathrm{two}}=\lceil T/n\rceil\) padded batches.  Let \(V_i\) be the
state after \(i\) batches.  At a batch the incoming \(A\) is
independent of \((S,V_i)\), and \(B\) is independent of the
whole actual experiment.  Apply Proposition~\ref{rlevels:two-comparison}
to the outgoing state \(V_{i+1}\); the required entropy bound follows
from \eqref{rlevels:clock} on the branch \(T<dL\).
Given \((D,DS)\), the exact batch \((A,AS)\) is known, so
the update is conditionally a channel from \(V_i\) without further
access to \(S\).  Therefore
\[
 \mathcal J(V_{i+1})
 \le I_{\mathrm{act}}(S;V_{i+1}\mid D,DS)+Cd
 \le I_{\mathrm{act}}(S;V_i\mid D,DS)+Cd
 =\mathcal J(V_i)+Cd.
\]
The last equality uses independence of \(D\) from \((S,V_i)\).
Set \(V_0=\overline U_0\).  It is independent of \(S\), and its
auxiliary \(G\) is independent of the pair \((S,V_0)\); hence
\(\mathcal J(V_0)=0\).  Iteration gives
\begin{equation}\label{rlevels:two-recurrence}
 \mathcal J(V_{B_{\mathrm{two}}})\le CdB_{\mathrm{two}}\le C'(T+d).
\end{equation}

For the independent \(k\)-row matrix defining \(\mathcal J\),
\eqref{eq:r-residual-markov} at \(\rho=1/4\) gives
\[
 \Pr\{R<1/4\}\le\frac{16}{15}\frac{k}{d}\le\frac8{15}.
\]
Apply Lemma~\ref{lem:r-residual-terminal} with \(p=2/3\) and
\(\eta=8/15\).  It gives
\begin{equation}\label{rlevels:two-endpoint}
 \mathcal J(V_{B_{\mathrm{two}}})
 \ge\frac2{15}(d-k-1)(L-\log8)-\log2.
\end{equation}
This is at least \(cdL\) uniformly for \(L\ge\log10\), since
\(L-\log8\ge(1-\log8/\log10)L>0\) and
\(d-k-1\ge d/3\).
The upper bound
\eqref{rlevels:two-recurrence} is \(O(T)\) by
\(T>d/4\).  This proves \(T\ge cdL\); the discarded branch
\(T\ge dL\) already has the same order.

In both arguments the auxiliary rows are analysis variables.  The
learner receives only its prescribed fresh Gaussian stream with exact
labels, and the comparisons permit arbitrary computation within a
transition subject to the persistent finite-state restriction.

\subsection{The hybrid comparison: stopping and terminal information}
\label{subsec:r-hybrid-application}

The averaged block estimate combines with the common stopping
reduction.  We need a terminal information bound uniform over the
auxiliary row counts in that average.

\Needspace{6\baselineskip}
\begin{lemma}[Information in a successful terminal state]
\label{rhyb:terminal}
Let \(V\) be a padded final state whose unit-vector output has
uniform-prior angular success at least \(2/3\) at
\(0<\epsilon\le1/10\).  For every \(k_0\le k\le k_1\),
\begin{equation}\label{rhyb:terminal-information}
 J_k(V)\ge\frac{d-k-1}{14}\log\frac1{5\epsilon}-\log2.
\end{equation}
\end{lemma}

\begin{proof}
For every \(k_0\le k\le k_1\), the auxiliary matrix has
\(k\le d/2\) rows.  At \(\rho=2/5\),
\eqref{eq:r-residual-markov} therefore gives
\[
 \Pr\{R<2/5\}\le\frac{k/d}{1-4/25}\le\frac{25}{42}.
\]
The success probability is at least \(2/3\), and
\(2/3-25/42=1/14\).  Lemma~\ref{lem:r-residual-terminal}
now gives \eqref{rhyb:terminal-information}, because
\(\rho/(2\epsilon)=1/(5\epsilon)\).  Its radius condition
holds for \(\epsilon\le1/10\).
\end{proof}

\begin{theorem}[Precision lower bound from hybrid projections]
\label{rhyb:conclusion}
Let \(M(d)\) be a nonnegative integer sequence with \(M(d)=o(d^2)\), and
let \(0<\epsilon(d)\le1/10\).  For each \(d\), consider a learner in
Definition~\ref{def:r-streaming-model} with memory bound \(M(d)\),
accuracy \(\epsilon(d)\), deterministic horizon \(T(d)\ge0\), and
uniform prior on \(S^{d-1}\).  If
\[
 \mathbb P\{\angle(\widehat S,S)\le\epsilon(d)\}\ge\frac23,
\]
then
\[
 T(d)\ge c\,d\log\frac1{\epsilon(d)}
\]
for all sufficiently large \(d\).  The constant \(c>0\) is absolute;
the dimension threshold may depend on the memory sequence.
\end{theorem}

\begin{proof}
Apply Lemma~\ref{lem:r-seed-reduction} with \(p_0=2/3\) to the
original jointly measurable experiment.  Fix a seed satisfying
\eqref{eq:r-seed-factorization} with conditional uniform-prior success
at least \(2/3\), retain its conditional initial-state law, and then
apply Lemma~\ref{lem:r-borel-kernels}.  The fixed-seed joint experiment,
state bound, and horizon are preserved.

Put \(b=\log(1/\epsilon)\).  In dimensions with \(T>db\), the conclusion
already holds.  In the remaining dimensions, Lemma~\ref{lem:r-stopping-reduction}
permits padding with \(W=(T+2)2^M\) states and
\(B=\log W=o(d^2)\), so eventually \(B\le d^2\).
Break the padded computation into blocks of \(h_0\) observations, with
a shorter final block if necessary.  The initial padded state
\(\overline U_0\) is independent of \(S\), and the side matrix is
independent of the pair \((S,\overline U_0)\).  Thus
\(J_k(\overline U_0)=0\) for every \(k\), and its averaged potential is zero.
Proposition~\ref{rhyb:block} gives
\[
 \overline J(V_{\rm final})
 \le Cd\left\lceil\frac{T}{h_0}\right\rceil
 \le C'(T+d)\le C''T.
\]
The last inequality uses the independent bound \(T>d/4\) from
Lemma~\ref{lem:r-stopping-reduction}; thus it also covers a shorter
last block.

Lemma~\ref{rhyb:terminal} gives the opposite bound for every \(k\) in
the average.  Since \(k\le d/2\) and \(b\ge\log10\),
\[
 \log\frac1{5\epsilon}=b-\log5
 \ge\frac{\log2}{\log10}\,b.
\]
The right side of \eqref{rhyb:terminal-information} is consequently at
least \(c'db\) for all sufficiently large \(d\), uniformly over the
allowed accuracies.  Comparing the two bounds proves \(T\ge cdb\).
The dimensions with \(T>db\) were already covered.
\end{proof}

If a learner succeeds for every unit signal with probability at least
\(2/3\), it also satisfies the uniform-prior hypothesis, so the theorem
applies to that worst-signal guarantee.  The shared-seed case is covered
by the joint-independence reduction at the start of the proof, which
retains fresh randomized transition and output kernels.  The proof uses
exact labels and imposes no linearity or within-transition precision
restriction.

\subsection{The finite Gaussian-fiber comparison}
\label{subsec:r-fiber-application}

We apply Proposition~\ref{rfib:block} to
Definition~\ref{def:r-streaming-model}, including its common notation
\(\Xi,U_0,(X_t)_{t\ge1}\).  This route also fixes the remaining
randomness after choosing the seed, so its terminal count uses a
deterministic decoder.

\begin{corollary}[Uniform-sphere sample lower bound]\label{rfib:streaming}
Let \(M(d)\ge0\) be integer-valued with \(M(d)=o(d^2)\), and let
\(0<\epsilon(d)\le1/10\).  Suppose learners in Definition~\ref{def:r-streaming-model},
with \(S\sim\sigma\), satisfy
\[
 \Pr\{\arccos\langle\widehat S,S\rangle\le\epsilon(d)\}\ge2/3,
\]
where the probability includes the signal, rows, and learner randomness.
Then
\[
 T(d)\ge c\,d\log(1/\epsilon(d))
\]
for an absolute \(c>0\) and all sufficiently large \(d\).  The
eventual dimension threshold may depend on the memory sequence.
\end{corollary}

\begin{proof}
Apply Lemma~\ref{lem:r-seed-reduction} with \(p_0=2/3\) to the
original jointly measurable uniform-prior experiment.  Fix a seed
\(\xi\) satisfying \eqref{eq:r-seed-factorization} whose conditional
success is at least \(2/3\), and retain the conditional initial-state
law.  Then apply Lemma~\ref{lem:r-borel-kernels} in that fixed-seed
experiment.  In particular all global Borel routing extensions are fixed
before \eqref{rfib:routing} and the finite measure \(Q\) are formed.
The two-point identity applies to those Borel functions without transferring
a null set from the real observation law to a perpendicular row law.

On a product extension of this fixed-seed experiment, choose an array
\(\mathcal R\) of independent uniform random variables consisting of
\(R_t\), \(0\le t\le T\), and
\(R^{\mathrm{out}}_{t,j}\), \(0\le t\le T\),
\(1\le j\le2^M\).  Choose the whole array independently of
\((U_0,S,(X_t)_{t\ge1})\).  The factorization
\eqref{eq:r-seed-factorization} then implies that the complete
conditional choice vector \((U_0,\mathcal R)\) is independent of
\((S,(X_t)_{t\ge1})\).

Enumerate each finite state alphabet by slots \(j\le2^M\).  Use
\(R_0\) for a stopping choice before the first sample and, at time
\(t\ge1\), use \(R_t\) to realize the joint next-state and stopping
outcome by the inverse distribution function of its finite probability
row.  Use \(R^{\mathrm{out}}_{t,j}\) to realize the unit-vector output
kernel at terminal pair \((t,j)\).  Such a realization can be chosen
Borel for a probability kernel on the standard Borel sphere, by mapping
the sphere into a Borel subset of \([0,1]\) and using an inverse
distribution function.  The resulting fixed-seed experiment is jointly
measurable and has the same conditional law as the Borel learner.

Its success probability averaged over \((U_0,\mathcal R)\) is at
least \(2/3\).  Fubini and averaging give a complete realization with
uniform-prior success at least \(3/5\).  Fix it.  The initial state,
all transition and stopping rules, and all terminal outputs are now
deterministic.  Independence of the complete choice vector leaves the
joint law of the signal and all rows unchanged.  Thus the complete vector
\((\xi,U_0,\mathcal R)\), including the seed and initialization, is
fixed before any state comparison.  Its unused entries add no persistent
learner state.

Put \(X_*=d\log(1/\epsilon)\), and restrict to \(T<X_*\).
Apply Lemma~\ref{lem:r-stopping-reduction} to the fixed decoder with
success at least \(3/5\).  Its output count gives
\begin{equation}\label{rfib:output-count}
 (d-1)\log(1/\epsilon)
 \le M\log2+\log(T+1)+\log(5/3).
\end{equation}
The common short-run estimate yields \(X_*=O(M+1)\), and the
absorbing padded state has width at most
\begin{equation}\label{rfib:width}
 \begin{aligned}
 N&=(T+2)2^M,\\
 \log N&\le M\log2+\log(X_*+2)=O(M+1)\le d^2
 \end{aligned}
\end{equation}
for all sufficiently large \(d\).  The padded decoder is a function
of this state, and rows generated after stopping remain fresh.

Pad also to \(\lceil T/n\rceil\) full \(n\)-row blocks.  Each block's
row matrix is independent of the pair consisting of the signal and old
state, while its labels remain exact.  After the complete fixation, the
initial padded state is deterministic, so its potential is zero.  Iterating
Proposition~\ref{rfib:block} for the padded final state \(F\) gives
\begin{equation}\label{rfib:iteration}
 \Phi(F)\le Kd\lceil T/n\rceil.
\end{equation}
For a lower bound use the independent \(k\)-row matrix in
\(\Phi(F)\).  Equation~\eqref{eq:r-residual-markov} gives
\[
 \Pr\{R<1/2\}\le\frac{4k}{3d}\le\frac49.
\]
The fixed decoder has success at least \(3/5\).  Apply
Lemma~\ref{lem:r-residual-terminal} with \(p=3/5\),
\(\rho=1/2\), and \(\eta=4/9\), and write
\(p_0=3/5-4/9>0\).  It gives
\begin{align}
 \Phi(F)
 &\ge p_0(d-k-1)\log(1/(4\epsilon))-\log2 \notag\\
 &\ge c\,d\log(1/\epsilon)                              \label{rfib:terminal}
\end{align}
for sufficiently large \(d\).  The constant is absolute because
\(d-k-1\ge d/2\) and
\(\log(1/(4\epsilon))\ge
 (1-\log4/\log10)\log(1/\epsilon)\) on the stated accuracy range.

Finally, Lemma~\ref{lem:r-stopping-reduction} gives
\(T>d/4\ge n\), including the exclusion of \(T=0\).
Hence \(\lceil T/n\rceil\le2T/n\), and \(d/n\) is
absolutely bounded for large \(d\).  The upper bound
\eqref{rfib:iteration} is therefore \(O(T)\).  Combining it with
\eqref{rfib:terminal} proves the corollary.
\end{proof}

If a learner has angular success at least \(2/3\) for every fixed
signal, averaging that guarantee over the uniform sphere supplies the
hypothesis of Corollary~\ref{rfib:streaming}.  It therefore gives the
same lower bound with the stated worst-signal premise.

